\documentclass[11pt]{article}
\usepackage[utf8]{inputenc}
\usepackage{float}
\usepackage{amsmath}
\newcommand{\mathbbm}[1]{\text{\usefont{U}{bbm}{m}{n}#1}} % from mathbbm.sty
\usepackage{amssymb}
\usepackage{soul} %to hilghlight with \hl{} and to strikeout with \st{}
\usepackage{color}  %to hilghlight with color with \hl{}
\usepackage{mathtools}
\usepackage{graphicx}
\usepackage{indentfirst}
\usepackage{enumitem}
\usepackage{setspace} 
%\setstretch{1.5}
\onehalfspacing
%\doublespacing
\usepackage{float}
\usepackage{fullpage}
\usepackage{ragged2e}
\usepackage[official]{eurosym}
\usepackage{booktabs}
\usepackage[bottom]{footmisc}
\newcommand{\sym}[1]{{#1}} % for symbols in Table
\usepackage{comment}
\usepackage{parskip}
\usepackage{tikz}
\usepackage[T1]{fontenc}
\usepackage{enumitem}
\setlist[itemize]{noitemsep}
\usepackage{placeins}
\usepackage[left=1in,right=1in,top=1in,bottom=1in]{geometry}


\usepackage{natbib}
\newcommand{\customcite}[1]{\citeauthor{#1}, \citeyear{#1}}
\bibliographystyle{chicago}
\usepackage{titlesec}
\usepackage[font=scriptsize]{subcaption}
	\captionsetup[figure]{font={small},labelfont={sc}}
	\captionsetup[table]{font={small},labelfont={sc}}
\titleformat*{\section}{\Large\bfseries}
\setcitestyle{authoryear,open={(},close={)}}
\usepackage{xcolor}
\usepackage[colorlinks]{hyperref}
\hypersetup{allcolors=blue!80!white}
\usepackage[scaled]{helvet}
\renewcommand{\familydefault}{\sfdefault}
 
\begin{document}
 
\begin{center}
    \textbf{\large \scshape CRISTIAN ARCIDIACONO}
\end{center}

% TODO: consider adding your office address here, e.g.
% Department of Economics, University of Bern, Bern, Switzerland
\begin{center}
    {\href{mailto:cristian.arcidiacono@unibe.ch}{cristian.arcidiacono@unibe.ch} $\vert$ \href{https://sites.google.com/view/cristian-arcidiacono/home}{Personal Page}}
\end{center}
\vspace{-0.2cm}
\noindent\rule{\textwidth}{0.4pt}

\vspace{0.2cm}

\begin{center}
    {\large \scshape Statement on Contributions to UCLA's Mission}
\end{center}

\vspace{0.3cm}

As a first-generation student who financed my own studies, I know from experience what the University of California's commitment to social mobility means. Before my Ph.D., I attended the Military Academy of Modena and then, while completing my bachelor's degree, ran the customer care department of a Milanese start-up for three years; that job funded my master's in Geneva. These experiences shape how I teach, how I do research, and how I serve the institutions I belong to.

\textbf{Teaching.} For nearly three years I have taught exercise sessions in ``International Finance'' and ``Money and Finance'' at UniDistance Suisse, a distance-learning university, and supervised bachelor's theses. Many of my students work alongside their studies or entered university through non-traditional paths. Without a physical classroom, I learned to plan each session carefully, prepare materials students can revisit, and check their understanding regularly. I would bring these practices to UCLA's classrooms, where many students share similar backgrounds, and I am committed to mentoring them as they find their own paths.

\textbf{Scholarship.} My research asks how scarcity in individual markets interacts with conditions shared across the economy: how labor shortages drive inflation, and how government borrowing affects the price of safe assets. These questions matter for households, whose wages and borrowing costs depend on the answers. Los Angeles, with one of the largest and most diverse labor markets in the country, is a natural setting in which to extend this work.

\textbf{Global reach and service.} Over the past four years I have worked in six cities and four countries, including research positions at the Bank for International Settlements and the European Stability Mechanism and a year at the LSE. These institutions taught me how a question changes when someone must act on its answer, and they gave me a network of economists in academia and policy that I would bring to UCLA's research community.

\textbf{Institutional effectiveness.} Running a team at a start-up taught me ownership, working under pressure, and turning ideas into results. I would carry that sense of responsibility into departmental service, from committee work to the mentoring of graduate students.

\vspace{0.3cm}

Cristian Arcidiacono

\end{document}
